\chapter{Нейроны по числу дендритов}
\chaptermark{Число дендритов}
\label{sec:dendrites}

\section{Число входов как параметр схемы}

В главе~\ref{sec:neurontypes} мы сортировали нейроны по тому, что выбрасывает их выход, а в главе~\ref{sec:eetypes} --- по тому, каким прибором они работают. Есть и третий признак, который инженер заметит сразу: \emph{сколько у нейрона входов}. Входы собирают дендриты --- ветви, на которых садятся чужие аксоны (глава~\ref{sec:dofa}, раздел~\ref{sec:weightstore}). У одних нейронов дендритов нет вовсе, у других один или два, у третьих --- дерево, собирающее сотню тысяч входов. Для электронщика это нагрузочная способность по входу, fan-in, и она почти всегда соответствует задаче.

Почему число и размер дендритов так важны, видно из простой оценки. Мембрана --- конденсатор с удельной ёмкостью $c_m\approx1$~мкФ/см$^2$, и чем больше площадь $A$ нейрона вместе с его дендритами, тем больше заряд нужно принести, чтобы поднять напряжение на $\Delta V$. Время, за которое входной ток $I$ это сделает, если не считать утечки,
\begin{empheq}[box=\keybox]{equation}
t \;\approx\; \frac{c_m\,A\,\Delta V}{I} .
\label{eq:dendriteload}
\end{empheq}
Маленький нейрон с несколькими короткими дендритами имеет площадь порядка $2\cdot10^{-5}$~см$^2$ и ёмкость около $20$~пФ. Один крупный синапс с током в несколько наноампер поднимает его на $20\mV$ меньше чем за десятую долю миллисекунды. У нейрона с большим деревом площадь раз в пятнадцать больше и ёмкость около $300$~пФ; обычный синапс с током около $0.1$~нА зарядил бы его за $60\ms$ --- много дольше постоянной времени утечки, то есть не зарядил бы никогда. Такому нейрону нужны десятки одновременных входов. Числа здесь --- порядки величин, но вывод от них не зависит: \emph{мало входов --- быстро и точно, много входов --- медленно и по сумме}.

\begin{figure}[t]
\centering
\begin{tikzpicture}[>=Stealth,
  soma/.style={circle,draw,very thick,fill=blue!6,minimum size=0.55cm,inner sep=0pt},
  ax/.style={->,very thick,red!70!black},
  den/.style={very thick,blue!60!black}]
% 0 dendrites: sensor on a line
\begin{scope}[xshift=0cm]
\draw[den,decorate,decoration={snake,amplitude=1pt,segment length=4pt}] (-0.9,0) -- (0,0);
\node[font=\scriptsize] at (-0.9,0.3) {датчик};
\draw[ax] (0,0) -- (1.0,0);
\draw[thick] (0.1,0) -- (0.1,0.45);
\node[soma,minimum size=0.4cm] at (0.1,0.65) {};
\node[font=\scriptsize,align=center] at (0.05,-0.75) {$0$\\линия с датчиком};
\end{scope}
% 1 dendrite
\begin{scope}[xshift=3.3cm]
\draw[den] (-0.9,0) -- (-0.28,0);
\node[soma] at (0,0) {};
\draw[ax] (0.28,0) -- (0.9,0);
\node[font=\scriptsize,align=center] at (0,-0.75) {$1$\\буфер, повторитель};
\end{scope}
% 2 dendrites: one per ear
\begin{scope}[xshift=6.2cm]
\draw[den] (-0.28,0) -- (-0.9,0); \draw[den] (0.28,0) -- (0.9,0);
\node[soma] at (0,0) {};
\draw[ax] (0,-0.28) -- (0,-0.6);
\node[font=\scriptsize] at (-0.9,0.25) {лев.}; \node[font=\scriptsize] at (0.9,0.25) {прав.};
\node[font=\scriptsize,align=center] at (0,-1.0) {$2$\\«И» по времени};
\end{scope}
% ~4 short dendrites: mixer
\begin{scope}[xshift=9.0cm]
\foreach \a in {45,135,225,315}{\draw[den] (\a:0.28) -- (\a:0.7); \fill[blue!60!black] (\a:0.72) circle (0.06);}
\node[soma] at (0,0) {};
\draw[ax] (0.28,0) -- (1.0,0);
\node[font=\scriptsize,align=center] at (0,-1.0) {$\sim4$\\смеситель};
\end{scope}
% many: summing tree
\begin{scope}[xshift=12.0cm]
\foreach \a in {60,90,120}{\draw[den] (0,0.28) -- ++(\a:0.55) coordinate (b); \foreach \d in {-20,20}{\draw[den] (b) -- ++(\a+\d:0.35);}}
\foreach \a in {-120,-60}{\draw[den] (0,-0.28) -- ++(\a:0.45);}
\node[soma] at (0,0) {};
\draw[ax] (0.28,0) -- (1.0,0);
\node[font=\scriptsize,align=center] at (0,-1.0) {$10^4$--$10^5$\\сумматор};
\end{scope}
\end{tikzpicture}
\caption{Пять классов по числу входов. Синим --- дендриты (входы), красным --- аксон (выход). Чем меньше входов, тем быстрее и точнее нейрон передаёт отдельный сигнал; чем больше, тем лучше он складывает и классифицирует. Схема, не рисунок формы клеток.}
\label{fig:dendriteclasses}
\end{figure}

\section{Ноль дендритов: датчик на конце линии}

У чувствительных нейронов осязания, боли и растяжения мышц (главы~\ref{sec:sensors}, \ref{sec:pain}, \ref{sec:muscles}) на теле клетки нет ни одного дендрита. Аксон выходит из тела одной веточкой и тут же делится надвое: одна ветвь уходит к коже или мышце, и её конец сам служит датчиком; другая идёт в спинной мозг. Импульс бежит от датчика прямо в спинной мозг, минуя тело клетки. Для инженера это линия связи, у которой датчик стоит на дальнем конце, а тело клетки висит сбоку, как блок питания: оно снабжает линию, но в передаче не участвует. Выгода --- скорость и надёжность: на пути нет ни одного места, где сигнал надо было бы суммировать и снова решать, передавать ли его.

Датчики света и звука --- ещё более крайний случай: это не нейроны-сборщики, а сами преобразователи, у которых входом служит не синапс, а белок-рецептор (рис.~\ref{fig:transducer}).

\section{Один дендрит: буфер}

Нейрон с одним дендритом и одним аксоном принимает одну входную линию и передаёт её дальше. Так устроены обонятельные нейроны: единственный дендрит выходит к поверхности носа и несёт единственный тип рецептора~\cite{Mombaerts2004}; так устроены передаточные клетки сетчатки между датчиками света и выходными нейронами глаза; так устроены нейроны, связывающие датчики уха с мозгом. Инженер назвал бы их буфером или повторителем: они не вычисляют, а доставляют один сигнал, иногда меняя его форму --- например, превращая плавное напряжение датчика в импульсы, как в главе~\ref{sec:sensors}.

\section{Два дендрита: «И» по времени}

Детекторы совпадений, определяющие направление на звук (рис.~\ref{fig:itd}), у млекопитающих часто имеют ровно два дендрита, направленных в противоположные стороны: на один приходят входы от левого уха, на другой --- от правого~\cite{Grothe2010}. Зачем разносить входы? Вспомним формулу~\eqref{eq:shunt}: когда на одном участке мембраны открыто много возбуждающих каналов, напряжение там приближается к равновесному потенциалу, и каждый следующий вход добавляет всё меньше. Поэтому много входов от одного уха, собранных на одном дендрите, насыщают его и не могут в одиночку довести нейрон до порога. А по одному входу с каждого дендрита складываются почти линейно. Два дендрита превращают нейрон в более строгое «И»: срабатывай, только если сигналы пришли с обеих сторон и вместе. На моделях показано, что такое разнесение действительно улучшает детекцию совпадений~\cite{AgmonSnir1998}.

\section{Несколько коротких дендритов: смеситель и ретранслятор}

\emph{Смесители.} В цепи обучения движениям из главы~\ref{sec:motorlearning} около семидесяти миллиардов маленьких \nt{ГЛУТ}-нейронов расширяют вход в очень длинный вектор. У каждого из них всего около четырёх коротких дендритов, и каждый заканчивается «когтем», который охватывает окончание одного входа. Каждый смеситель, таким образом, видит лишь около четырёх входов из многих и работает как маленький элемент «И» над своей четвёркой. Из $M$ входных линий можно выбрать $\binom{M}{4}$ разных четвёрок: при $M=100$ это почти четыре миллиона. Зачем столько? Пороговый элемент с $n$ входами может правильно разделить на два класса не больше
\begin{empheq}[box=\keybox]{equation}
P_{\max} \;\approx\; 2n
\label{eq:cover}
\end{empheq}
случайных образов~\cite{Cover1965}. Выходной нейрон той же цепи получает около $10^5$ входов от смесителей и может выучить порядка $2\cdot10^5$ разных соответствий. Смесители с малым числом входов нужны именно для того, чтобы у большого сумматора было много разных, почти независимых входов~\cite{Marr1969,Albus1971}.

\emph{Ретрансляторы.} На пути от уха к детекторам направления есть нейроны с немногими короткими дендритами, которые получают всего несколько огромных синапсов, а некоторые --- по существу один. По~\eqref{eq:dendriteload} это ровно то, что нужно: маленькая ёмкость и большой ток, и нейрон срабатывает через доли миллисекунды после каждого входного импульса, почти не теряя точности времени. Один из таких ретрансляторов получает \nt{ГЛУТ}, а выдаёт \nt{ГЛИЦ}: это инвертор с точностью в десятки микросекунд, поставляющий быстрое торможение детекторам направления~\cite{BorstSoria2012}.

\section{Тысячи входов: сумматор и классификатор}

На другом конце шкалы --- \nt{ГЛУТ}-нейроны коры с десятком тысяч входов и выходные \nt{ГАМК}-нейроны цепи обучения движениям с сотней тысяч. Такой нейрон по~\eqref{eq:dendriteload} медленный и не может откликнуться на отдельный вход: он складывает. Это интегрирующий нейрон главы~\ref{sec:glutcomputer} и тот сумматор, из которого строятся классификаторы. Большое дерево добавляет и новое: его ветви работают как отдельные подсхемы со своими порогами, так что один нейрон ведёт себя как маленькая многослойная сеть~\cite{Beniaguev2021,Gidon2020}.

\section{Дендриты, которые ещё и выход}

Есть нейроны, у которых аксона нет вовсе, а дендриты служат и входом, и выходом: они получают сигналы и сами выбрасывают медиатор. Самый изученный пример --- \nt{ГАМК}-нейроны сетчатки, участвующие в определении направления движения (глава~\ref{sec:gaba}, рис.~\ref{fig:direction}). У такого нейрона каждая ветвь сама по себе отвечает на движение от тела клетки к своему кончику и выдаёт торможение с этого кончика~\cite{Euler2002}. Один нейрон --- это несколько независимых детекторов направления, по одному на ветвь. Для инженера это не один элемент, а микросхема с несколькими каналами в одном корпусе.

\section{Итог}

\begin{table}[t]
\centering
\footnotesize
\setlength{\tabcolsep}{3pt}
\renewcommand{\arraystretch}{1.15}
\begin{tabular}{>{\raggedright}p{1.9cm}>{\raggedright}p{3.4cm}>{\raggedright}p{2.6cm}>{\raggedright}p{1.8cm}>{\raggedright\arraybackslash}p{3.8cm}}
\toprule
Дендритов & Пример & Прибор & Входов & Что известно \\
\midrule
$0$ & датчики осязания, боли, растяжения & линия с датчиком на конце & --- & установлено \\
$1$ & обонятельные нейроны, передаточные клетки сетчатки, нейроны уха & буфер, повторитель & $1$--немного & установлено \\
$2$ & детекторы направления на звук & «И» по времени & два пучка & строение установлено; выигрыш от разнесения входов показан на моделях \\
$\sim4$ & смесители цепи обучения движениям & маленький элемент «И» & $\sim4$ & установлено; роль расширения входа --- хорошо обоснованная гипотеза \\
немного, короткие & ретрансляторы на пути от уха & ретранслятор, инвертор & $1$--несколько огромных & установлено \\
тысячи & \nt{ГЛУТ}-нейроны коры, выходные нейроны обучения движениям & сумматор, классификатор & $10^4$--$10^5$ & установлено; ветви как подсхемы измерены на отдельных примерах \\
дендриты-выходы & \nt{ГАМК}-нейроны направления в сетчатке & многоканальная микросхема & много & измерено \\
\bottomrule
\end{tabular}
\caption{Нейроны по числу дендритов.}
\label{tab:dendrites}
\end{table}

\begin{keyblock}
\begin{proposition}[Число входов следует за задачей]
\label{prop:dendrites}
Где нужна скорость и точность отдельного сигнала --- на датчиках, ретрансляторах и детекторах совпадений, --- у нейронов мало дендритов, мало входов и большие синапсы: малая ёмкость~\eqref{eq:dendriteload} заряжается быстро. Где нужно складывать и классифицировать, у нейронов большие деревья с тысячами входов. Строение классов измерено; вычислительное объяснение хорошо обосновано, но для многих классов остаётся гипотезой.
\end{proposition}
\end{keyblock}

Таблица~\ref{tab:dendrites} дополняет две прежние сортировки: по медиатору (таблица~\ref{tab:types}) и по прибору (таблица~\ref{tab:eetypes}). Все три смотрят на один и тот же элемент с разных сторон: что он выдаёт, что он вычисляет и сколько он слушает.

\section{Задачи}

\begin{exercise}[\lvlA]
\label{ex:19:1}
\emph{Сколько входов нужно большому нейрону.} Нейрон с большим деревом имеет ёмкость $C=300$~пФ и постоянную времени утечки $\tau_m=20\ms$; каждый обычный синапс даёт ток $0.1$~нА (считайте синапсы источниками тока). Порог --- $20\mV$ над покоем.
(а)~Найдите сопротивление мембраны и установившийся сдвиг напряжения от одного синапса.
(б)~Сколько одновременно работающих синапсов нужно, чтобы дойти до порога вообще; за $10\ms$; за $5\ms$?
(в)~Маленький нейрон ($C=20$~пФ, тот же $\tau_m$) получает один синапс в $4$~нА. За какое время он дойдёт до порога и почему здесь утечкой можно пренебречь?
\end{exercise}

\begin{exercise}[\lvlA]
\label{ex:19:2}
\emph{Бюджет смесителей.} (а)~Проверьте, что из $M=100$ входных линий можно выбрать $\binom{100}{4}\approx3.9\cdot10^6$ четвёрок.
(б)~Выходной нейрон с $n=10^5$ входами по~\eqref{eq:cover} запоминает $P_{\max}$ соответствий по одному биту. Сколько бит приходится на один синапс?
(в)~Во сколько раз расширение входа через смесители увеличивает число соответствий, которые может выучить выходной нейрон, по сравнению с нейроном, получающим прямо $M=100$ исходных линий?
\end{exercise}

\begin{exercise}[\lvlB]
\label{ex:19:3}
\emph{Откуда берётся $2n$.} Число способов разбить $P$ точек общего положения в $n$-мерном пространстве на два класса гиперплоскостью, проходящей через начало координат, равно
\begin{equation*}
C(P,n)=2\sum_{k=0}^{n-1}\binom{P-1}{k}.
\end{equation*}
(а)~Покажите, что при $P\le n$ разделимы все $2^P$ разбиений.
(б)~Докажите, что при $P=2n$ разделима ровно половина разбиений.
(в)~Приблизив биномиальное распределение нормальным, найдите долю разделимых разбиений при $n=100$ и $P=180$, $220$. Какова ширина перехода от «почти все» к «почти никакие» и как она зависит от $n$? Почему~\eqref{eq:cover} становится точным порогом при больших $n$?
\end{exercise}

\begin{exercise}[\lvlB]
\label{ex:19:4}
\emph{Два дендрита как строгое «И».} Упростим детектор направления на звук. Каждый дендрит --- отдельный участок с проводимостью утечки $g_L$; синапсы открывают на нём возбуждающую проводимость, и напряжение участка равно~\eqref{eq:shunt} без торможения. Тело клетки складывает напряжения двух участков с коэффициентом связи $\kappa$: $V_{\text{тело}}=\kappa(V_1+V_2)$. Для сравнения точечный нейрон --- один участок, на который приходят все входы.
(а)~Каждый вход открывает проводимость $g_0$. Покажите, что точечный нейрон не отличает $2n$ входов с одной стороны от $n+n$ с двух сторон.
(б)~Найдите для нейрона с двумя дендритами отношение отклика на $n+n$ к отклику на $2n+0$ как функцию $G=ng_0/g_L$ и его предел.
(в)~Выберите порог $\theta=1.1\,\kappa E_E$. Покажите, что никакое число входов с одной стороны не доводит нейрон до порога, и найдите наименьшее $n$ на каждой стороне, при котором он срабатывает, если $g_0=0.5\,g_L$.
\end{exercise}

\begin{exercise}[\lvlB]
\label{ex:19:5}
\emph{Проектирование ретранслятора.} Дрожание момента срабатывания равно $\sigma_t=\sigma_V/(\dd V/\dd t)$, где $\sigma_V$ --- шум напряжения, а производная берётся в момент пересечения порога. Требуется ретранслятор с $\sigma_t\le20$~мкс при $\sigma_V=1\mV$.
(а)~Какой наименьший синаптический ток нужен нейрону с $C=20$~пФ? Утечку на интервале в доли миллисекунды не учитывайте.
(б)~Сколько обычных синапсов по $0.1$~нА должно сработать одновременно, чтобы той же точности достиг нейрон с $C=300$~пФ? Реалистично ли это?
(в)~Каково дрожание нейрона с $C=20$~пФ и одним синапсом в $4$~нА? Сопоставьте с «точностью в десятки микросекунд» из текста.
\end{exercise}

\begin{exercise}[\lvlB]
\label{ex:19:6}
\emph{Моделирование: заряд длинного дендрита.} Пассивный дендрит длины $\ell$ описывается уравнением кабеля
\begin{equation*}
\tau_m\frac{\partial V}{\partial t}=\lambda^2\frac{\partial^2V}{\partial x^2}-V+\frac{i(x,t)}{g},
\end{equation*}
где $\lambda$ --- постоянная длины, $g$ --- проводимость утечки на единицу длины; концы закрыты ($\partial V/\partial x=0$). Перейдите к безразмерным $T=t/\tau_m$, $X=x/\lambda$, $L=\ell/\lambda$.
(а)~Выведите установившееся отношение $V(0)/V(L)$ при постоянном токе, поданном в дальний конец.
(б)~Напишите на Python модель из $40$ участков (Эйлер, шаг $\Delta T\le0.2\,\Delta X^2$). Подайте в дальний конец короткий импульс тока длительностью $0.01\,\tau_m$ и измерьте на ближнем конце ($X=0$) время максимума и его величину в долях значения $Q/(c_m\ell)$, которое дал бы тот же заряд $Q$, мгновенно размазанный по всему дендриту. Сделайте это для $L=0.2$, $1$ и $2$ и проверьте результат (а).
(в)~Когда оценка~\eqref{eq:dendriteload} по полной ёмкости верна, а когда нет? Что это говорит о коротких дендритах ретрансляторов?
\end{exercise}

\begin{exercise}[\lvlC]
\label{ex:19:7}
\emph{Исследование: оптимальное число входов.} Утверждение~\ref{prop:dendrites} говорит, что вычислительное объяснение числа входов для многих классов остаётся гипотезой. Постройте количественную модель компромисса.
(а)~Пусть ёмкость нейрона $C=C_0+nc_s$ растёт с числом входов $n$, одновременно работает $m$ входов по току $I_s$, и нейрон должен запоминать $P\approx2n$ соответствий~\eqref{eq:cover}. Выразите время отклика через $P$ и найдите, что меняется, если работает не фиксированное число входов, а фиксированная доля $m=\varphi n$.
(б)~Предложите критерий стоимости (время, энергия заряда $CV^2$, ёмкость памяти, дрожание из задачи~\ref{ex:19:5}) и найдите оптимальное $n$ для ретранслятора, детектора совпадений и классификатора. Какие строки таблицы~\ref{tab:dendrites} модель объясняет, а какие нет?
(в)~Большое дерево с ветвями-подсхемами ведёт себя как маленькая многослойная сеть. Предложите численный опыт, сравнивающий ёмкость нейрона из $m$ ветвей по $k$ входов с пороговой нелинейностью в каждой ветви и ёмкость точечного нейрона с $n=mk$ входами.
\end{exercise}
